\subsection{Summary}
\label{Subsection:Summary}

Three things were built and one question was answered.

The first is a trading framework in which one strategy implementation runs in three places: an offline backtest over stored quotes, the cTrader backtesting tab, and a live account. Its execution engine reproduces the platform's own results byte for byte across six reference configurations and extends them, stepping on ticks rather than bar closes and charging the spread recorded at the moment of each trade. Behind it sits a market database of $1.66$ billion quotes across the seven pairs, occupying $295$ GB, from which every bar series used in this work is derived.

The second is an agent. It observes thirty-eight dimensionless values built from sixteen technical indicators and emits one continuous action that sets the direction and the size of its position together. Six deviations from the published algorithm were needed to make it train on this problem, of which two matter: a scale-sensitive reward, because a scale-invariant one gives no gradient towards holding a meaningful position, and a penalty on the actor's pre-activation, without which the policy saturates and collapses to a constant action.

The third is a selection procedure. Candidates are rejected by gates that remove degenerate policies, tested against a null built by rotating each model's own position sequence in time, and ranked by a composite in which return is one component of nine. The procedure was not built to be elegant; it was built because ranking on return does not work, which Section~\ref{Subsection:StatisticalScreening} demonstrates rather than asserts.

The question was answered as follows. All seven agents are profitable over the eleven-year window and from all five starting balances, with total returns from $+15.7\%$ to $+226.6\%$ net of the spread and commission of a swap-free raw-spread account. Six of the seven produce positive alpha once their return is decomposed against their own market, and three of those carry almost no market exposure while doing it. The seventh, USD/JPY, reproduces its market and is reported as a failure with a diagnosed cause. Combining the seven into an equal-weight portfolio reduces the maximum drawdown from an average of $34.5\%$ to $15.9\%$ and raises the Sharpe ratio to $0.632$, because the agents are close to uncorrelated with one another.

Two findings sit underneath those numbers and are more transferable than they are. The continuous action is genuinely used: the mean absolute action ranges from $0.09$ to $0.98$ across the pairs and saturation is below one percent everywhere, so the agents learned to size positions rather than to switch between extremes. And win rates lie between $49.8\%$ and $55.0\%$, with the strongest model in the set having the lowest of them, so the edge does not come from being right more often. It comes from what the agents do with a position once it is open, which is the capability a continuous action provides and a discrete one does not.